



We can create maps between projective spaces. 
There are two possible ways to do so.
The first ist by using the \verb'-map' command, which is a projective space activity.
The second is by using a map object.
We will demonstrate both ways next.


\bigskip


The first example shows a Cremona map from $\PG(2,13)$ to $\PG(3,13).$ 
At first, we define the 4 coordinate functions as makefile variables. 
These coordinate functions must be polynomials. 
In the command, we create the finite field, the projective space for the domain of the map, 
and the polynomial ring associated with the domain. 
After that, we will create symbolic objects for the 4 coordinate functions. 
The object called ``Cremona'' is a vector of the coordinate functions. 
Finally, the \verb'-map' command is used to define and evaluate the map. 
This command is a projective space activity.
Here are the necessary makefile definitions and commands:

\sourcecodevariable{CREMONA_MAP_Y0}


\sourcecodevariable{CREMONA_MAP_Y1}


\sourcecodevariable{CREMONA_MAP_Y2}


\sourcecodevariable{CREMONA_MAP_Y3}


\sourcecode{Cremona_map}



\bigskip






Next, we demonstrate the mapping object 
to define a map between projective spaces. 
Table~\ref{tab:mapping} lists Orbiter 
commands for mappings between projective spaces.
\orbitertableofcommands{mapping}


The first example creates the Veronese variety of $\PG(1,7)$ in $\PG(3,7)$:

\sourcecode{Veronese_1_3_q7}

The image is written as a set of points in Orbiter ranks:
\orbiteroutput{GRAPHICS/LATEX/Veronese_1_3_q7}
For a deeper exploration of the combinatorial and group-theoretic properties of this object, 
see Section~\ref{sec:canonical:form:objects:incidence:geometries}.

\bigskip

